% !TEX root = ../main.tex
\chapter{Conclusions i línies futures}
\label{chap:conclusions}

\section{Resposta a la pregunta de recerca}

La pregunta que vertebra el treball és aquesta: \emph{sota un mateix simulador, format i protocol, quin dels cinc paradigmes ---heurístiques, minimax d'un torn, MCTS ras, imitació i MaskablePPO--- obté millor rendiment observat a \texttt{gen9randombattle}, i què explica les diferències?}

Heurístiques, minimax, MCTS ras i imitació, amb sis agents de referència, formen una matriu dirigida de $28\times 28$. MaskablePPO s'avalua a part, contra \texttt{v12}. La comparació comparteix Showdown i el format, però no és simètrica: l'MCTS té $n=1.000$ per cel·la i el PPO no entra a la lliga.

Resposta: sota les implementacions i els pressupostos d'aquest experiment, les heurístiques obtenen el millor resultat observat. \texttt{v12} assoleix 69,01\% de mitjana global descriptiva i 59,91\% contra Abyssal. Contra el mateix comparador \texttt{v12}, \texttt{v17} (minimax) obté 40,94\%, \texttt{v20} (MCTS) 42,7\%, \texttt{v21} (imitació) 41,28\% i el model PPO 34,1\% ($n=10.000$ tret de l'MCTS).

Això no afirma que les heurístiques siguin universalment superiors. Afirma que, aquí, el coneixement explícit sobre Teracristal·lització, selecció del substitut, debilitaments garantits i control del ritme no el recuperen ni la cerca rasa ni els models apresos amb els recursos utilitzats.

\subsection{Grau d'acompliment dels objectius específics}

A partir dels resultats empírics obtinguts, es pot verificar el compliment dels sis objectius específics plantejats en la Introducció (Secció~\ref{sec:objectius}):

\begin{enumerate}
    \item \textbf{Formalització teòrica del domini (Acomplert):} S'ha formalitzat el combat Pokémon com un joc estocàstic parcialment observable de dos jugadors (POSG) al Capítol~\ref{chap:marc_teoric}, i s'ha justificat la reducció a POMDP quan un sol decisor absorbeix la política rival $\pi_2$ dins de $T$. S'han definit les fonts d'incertesa de l'estat ocult, la dinàmica d'accions simultànies i les transicions estocàstiques.
    \item \textbf{Disseny metodològic, infraestructura i telemetria (Acomplert):} S'ha dissenyat el protocol experimental i el model de rànquing Bradley--Terry (Capítol~\ref{chap:metodologia}), s'ha construït la infraestructura computacional paral·lela i recuperable (Capítol~\ref{chap:infraestructura}) que ha produït 6.409.000 batalles dirigides a la matriu, i s'ha instrumentat el marc de telemetria per torn per a l'explicabilitat conductual.
    \item \textbf{Paradigmes heurístics i de cerca en línia (Acomplert):} S'ha desenvolupat la seqüència d'ablacions heurístiques (\texttt{v1}--\texttt{v14}) i s'han implementat i instrumentat les variants de cerca minimax d'un torn (\texttt{v15}--\texttt{v17}) i MCTS ras (\texttt{v18}--\texttt{v20}) amb priors d'orientació (Capítol~\ref{chap:implementacio}). \texttt{v12} és la millor heurística observada; \texttt{v14} és l'heurística de servei (fulla, prior PUCT, executor d'imitació i bot del \emph{ladder}).
    \item \textbf{Agents apresos i integració de LLMs (Acomplert):} S'han entrenat els models d'imitació XGBoost (\texttt{v21}--\texttt{v22}), s'ha optimitzat MaskablePPO amb clonatge conductual sobre \texttt{v8} i s'ha avaluat la viabilitat exploratòria de \textit{PokéChamp} i \textit{PokéLLMon}.
    \item \textbf{Avaluació empírica massiva i rànquings (Acomplert):} S'ha completat l'avaluació de la matriu competitiva dirigida de 28 agents ($28\times 28$), calculant els rànquings Bradley--Terry i analitzant el comportament obtingut per telemetria (Capítol~\ref{chap:resultats}).
    \item \textbf{Avaluació externa al ladder públic (Acomplert):} S'ha desplegat l'heurística de servei (\texttt{v14}) a la classificació pública de Pokémon Showdown, amb 431 batalles contra jugadors humans (usuari \texttt{SirPThesis}). El desplegament no pretén que \texttt{v14} sigui l'agent de millor taxa; és la política més completa i instrumentada per operar en línia.
\end{enumerate}

\section{Resultats principals i avaluació de les hipòtesis}
\label{sec:resultats_principals_hipotesis}

A continuació es contrasten els resultats empírics amb les cinc hipòtesis de recerca formalitzades a la Secció~\ref{subsec:hipotesis}:

\begin{enumerate}
    \item \textbf{$\mathbf{H_1}$ (Monotonicitat de l'ablació heurística) --- Refutada: La complexitat heurística no produeix una millora monotònica.} \texttt{v1}--\texttt{v6} formen un altiplà d'efectivitat similar; \texttt{v7}, \texttt{v9} i \texttt{v12} introdueixen els salts de rendiment més rellevants. Tanmateix, les versions més complexes (\texttt{v13} i \texttt{v14}) afegeixen una càrrega de regles que, lluny de millorar el joc, redueix la taxa global (62,03\% per a \texttt{v14} enfront del 69,01\% de \texttt{v12}), compatible amb una especialització de domini (més inferència i model de rival) que no ajuda contra polítiques deterministes.
    \item \textbf{$\mathbf{H_2}$ (Cerca adversària en línia) --- Refutada: El minimax d'un torn i l'MCTS ras no superen l'heurística de partida.} Ni la matriu simultània minimax (\texttt{v15}--\texttt{v17}) ni les simulacions locals MCTS a cinc torns (\texttt{v18}--\texttt{v20}) aconsegueixen batre la política que avalua les fulles (\texttt{v14}) ni l'heurística òptima (\texttt{v12}). \texttt{v17} obté un 43,61\% contra \texttt{v14} i un 40,94\% contra \texttt{v12}, mentre que \texttt{v20} n'obté un 42,7\%. La incertesa de l'estat ocult rival i el soroll estocàstic limiten la utilitat de la cerca a poca profunditat.
    \item \textbf{$\mathbf{H_3}$ (Guiatge per priors a l'MCTS) --- Confirmada: L'MCTS ras depèn críticament del biaix heurístic.} La cerca sense informació prèvia (UCB1) discrepa de les decisions expertes de \texttt{v14} en el 71,3\% de les decisions de cerca (\texttt{v18}), generant branques subòptimes. La incorporació de priors heurístics mitjançant l'algorisme PUCT redueix aquesta taxa de discrepància dràsticament al 18,7\% i millora la puntuació Bradley--Terry de la família (1742 per a \texttt{v20} vs. 1697 per a \texttt{v18}).
    \item \textbf{$\mathbf{H_4}$ (Bastida de domini a la imitació) --- Confirmada: La imitació requereix regles expertes d'execució.} L'agent híbrid \texttt{v21} (XGBoost per a la macrodecisió de canviar i \texttt{v14} per a l'execució micro) supera l'agent d'imitació pura \texttt{v22} en tots i cadascun dels 28 enfrontaments de la matriu, registrant una taxa global del 58,71\% enfront del 33,49\%. Les dades humanes capturen el ritme global, però necessiten càlcul de dany determinista per a la microexecució.
    \item \textbf{$\mathbf{H_5}$ (Currículum i DRL) --- Parcialment refutada: El currículum PPO supera les etapes inicials però s'estanca davant d'heurístiques avançades.} El model MaskablePPO amb clonatge conductual sobre \texttt{v8} supera amb solvència les etapes 1 i 1.5 del currículum (contra Random i MaxBasePower), però s'estanca a la fase 2 contra \texttt{v8} (40,6\% de victòries, per sota del llindar de traspàs del 55\%). En el test final de la fase 3 contra \texttt{v12}, el model assoleix un 34,1\% ($n=10.000$), un valor molt proper al de l'agent de clonatge \texttt{v8} (32,9\%), indicant que l'RL no aconsegueix transcendir la política base sota els pressupostos d'entrenament avaluats.
    \item \textbf{Síntesi transversal: La integració residual és el patró més consistent.} Els agents de cerca i aprenentatge amb millor rendiment (\texttt{v17}, \texttt{v20} i \texttt{v21}) són precisament aquells que conserven una part substancial de coneixement heurístic com a política residual o de servei. En jocs POSG d'alta combinatòria i observabilitat parcial, el coneixement de domini expert actua com a condició de contorn indispensable.
\end{enumerate}

\section{Contribucions}

\begin{itemize}
    \item una infraestructura coordinador--processos de treball, recuperable i instrumentada, que ha produït 6.409.000 batalles a la matriu dirigida;
    \item catorze heurístiques que documenten la incorporació de coneixement específic del format, amb \texttt{v12} com a millor observació i \texttt{v14} com a política de servei;
    \item variants instrumentades de minimax d'un torn i MCTS ras, amb telemetria de freqüència i discrepància;
    \item dos agents d'imitació entrenats amb dades del mateix format i una anàlisi del paper de l'executor;
    \item un entorn d'accions emmascarades, clonatge conductual i MaskablePPO sobre Showdown;
    \item un protocol que separa l'avaluació local, l'experiment PPO i la mostra pública contra persones.
\end{itemize}

\section{Limitacions}

Els cinc paradigmes no participen en una mateixa lliga: MaskablePPO només s'avalua contra \texttt{v12} i l'MCTS utilitza una mostra menor ($n=1.000$ per cel·la). Les taxes globals tenen ponderacions diferents; no hi ha equips o llavors aparellats; es fan moltes comparacions exploratòries; i les millors variants es seleccionen després d'observar la matriu. Tot i que versions com \texttt{v12} introdueixen Teracristal·lització i selecció de substituts simultàniament, la telemetria comparativa amb \texttt{v13} permet descompondre observacionalment el pes de cadascuna; el llindar d'imitació es tria sobre la partició retinguda, i no s'ha repetit l'entrenament complet d'IL amb diverses llavors. La mostra del \emph{ladder} (431 partides, 39,44\% $\pm$ 4,6~pp) és insuficient per estimar un Elo estable. Finalment, els resultats només descriuen \texttt{gen9randombattle} i les versions de programari utilitzades: un POSG de torns simultanis, per equips, amb informació imperfecta. No es transfereixen automàticament a formats amb equips tancats ni a Dobles.

\section{Línies futures}

Les continuacions prioritàries s'organitzen en tres eixos que tanquen forats de l'experiment actual.

\subsection{Cerca, creences i priors}

La cerca d'aquest treball és rasa: profunditat~$d=1$ i \emph{rollouts} de cinc torns. No és un arbre adversari profund.

\begin{itemize}
    \item Substituir l'UCB1 només a l'arrel per un arbre de simulació complet o per \textit{Information Set MCTS} (ISMCTS), capaç de gestionar l'estat ocult del rival mitjançant mostratge de creences (\emph{determinization}).
    \item Aprendre polítiques de simulació (\emph{rollout policies}) o priors de fulla a partir de \texttt{v12}, la millor heurística observada, i no només de \texttt{v14}. \texttt{v14} continua sent el candidat natural com a executor i com a prior PUCT ja instrumentat; el millor resultat de l'escala, però, és \texttt{v12}.
    \item Mesurar el cost per decisió (temps de paret i energia) a profunditats creixents, per no confondre un millor algorisme amb un pressupost de cerca més gran.
\end{itemize}

\subsection{Aprenentatge per reforç i imitació}

\begin{itemize}
    \item Repetir l'entrenament MaskablePPO amb clonatge conductual sobre \texttt{v12} ---no sobre \texttt{v8}--- i amb xarxes recurrents (LSTM/GRU) que mantinguin memòria de l'estat ocult al llarg de l'episodi.
    \item Entrenar un model d'imitació condicionat al conjunt de candidats legals, capaç de decidir simultàniament moviments, canvis i Teracristal·lització sense un executor heurístic extern. Això aïllaria l'efecte que ara barreja \texttt{v21} (classificador de canvi) i \texttt{v14} (execució).
    \item Completar una matriu simètrica que inclogui PPO i models de llenguatge sota pressupostos de còmput equivalents, mesurant el cost temporal i financer per decisió. Sense aquesta lliga, la taxa del 34,1\% de PPO contra \texttt{v12} no és comparable amb les taxes globals de la matriu de 28.
\end{itemize}

\subsection{Avaluació: llavors i persones}

\begin{itemize}
    \item Utilitzar equips i llavors aleatòries aparellades (\emph{paired seed benchmarks}) per reduir la variabilitat de la generació procedural.
    \item Ampliar l'avaluació al \emph{ladder} públic amb diversos agents ---inclosa \texttt{v12}--- i finestres temporals més llargues, per obtenir rànquings Elo estables i no una mostra de 431 partides d'un sol bot.
\end{itemize}

% ============================================================
\section{Consideracions ètiques, socials i d'impacte ambiental}
\label{sec:etica_impacte}

Aquest treball toca tres qüestions que no es resolen amb el rànquing de la matriu: si era just posar un agent al \emph{ladder} públic, què queda d'útil fora de Random Battle, i quina energia ha calgut per mesurar-ho.

\subsection{Ètica i fair play en entorns de joc en línia}

Les 6.409.000 batalles de la matriu es van jugar entre agents, en servidors Showdown propis. No ocupen la infraestructura pública ni mouen l'Elo de ningú. El problema ètic, si n'hi ha un, és la mostra externa: 431 partides de \texttt{v14} contra persones, els dies 9 i 10 de juny de 2026, amb el compte \texttt{SirPThesis} (Seccions~\ref{sec:disseny_experiments} i~\ref{sec:deployment}).

Qui s'hi enfrontava no havia donat un consentiment informat. El bot pot alterar-li la puntuació i, si no llegeix el nom del compte, pot prendre'l per un jugador. Pokémon Showdown no té un registre per a bots de combat ---el rang \emph{Global Bot} és per a moderació de xat---, de manera que l'agent entra pel mateix WebSocket que qualsevol usuari. Això no és un permís en blanc.

S'ha parlat amb l'equip de Showdown i Smogon abans del desplegament. El compte indica que es tracta d'una tesi; no s'ha participat en tornejos; s'ha respectat el límit de 12 partides cada 3 minuts; i la plataforma ja exclou el tràfic de bots de les estadístiques d'ús del format. El nom redueix l'asimetria, però no l'elimina: no s'ha anunciat a cada sala que l'adversari era un programa. Per a un TFM, n'hi ha prou amb ser identificable i no interferir en la competició oficial. No n'hi ha prou per afirmar que un agent autònom al \emph{ladder} sigui inofensiu en general.

\subsection{Impacte social i transferibilitat científica}

L'impacte social directe és modest, i convé dir-ho sense envoltar-ho d'aplicacions inventades. Aquest treball no millora un servei públic. Afecta, de manera puntual, qui es va trobar \texttt{SirPThesis} a la cua i, de manera més duradora, qui vulgui reproduir o estendre el benchmark: el protocol, els 28 agents i la telemetria queden documentats.

El que es pot portar a un altre POSG de dos jugadors ---accions simultànies, observació parcial--- no és el catàleg de regles de la Generació~IX. És la manera de comparar: un mateix contracte d'entorn per a tots els paradigmes, comptadors per torn que distingeixen \emph{què} fa l'agent de \emph{si} guanya, i el patró residual de conservar restriccions de domini en lloc de substituir-les d'entrada per cerca rasa o per una política apresa. Les taxes 69,01\% i 34,1\% no viatgen: descriuen \texttt{gen9randombattle} sota els pressupostos d'aquest experiment.

\subsection{Impacte ambiental i petjada de carboni}

El còmput és el del Capítol~\ref{chap:infraestructura} (Secció~\ref{sec:cost_economic}): 685,04~hores de màquina, 135~W de mitjana i \textbf{92,5~kWh} en total, de les quals 471,38~h corresponen a la matriu de Generació~IX. Aproximadament la meitat de l'energia ha vingut de la xarxa, de nit (5,85~€); l'altra meitat, de l'autoconsum solar de l'habitatge. Amb un factor conservador de 0,21~kg $\mathrm{CO}_2$/kWh, la porció de xarxa ($\approx 46$~kWh) equival a uns \textbf{9,7~kg $\mathrm{CO}_2$e}. Si s'atribuís tot el 92,5~kWh al mix, el sostre seria 19,4~kg. Són xifres d'electricitat d'operació, no un cicle de vida del maquinari.

No s'ha repetit l'experiment en un clúster de GPU, així que no es pot parlar d'un percentatge d'estalvi. El que sí que pesa, ambientalment i metodològicament, és haver comparat els cinc paradigmes amb cerca rasa i un sol lloc de treball: el cost queda en desenes de quilowatts-hora, no en l'escala habitual de l'aprenentatge per reforç a gran escala.

% ============================================================
\section{Tancament}

El resultat central és metodològic i tècnic: comparar paradigmes exigeix fer explícits el coneixement incorporat, el pressupost de còmput, la representació de l'estat i el protocol d'avaluació. En aquest experiment, conservar coneixement específic de Random Battle és més efectiu que substituir-lo per cerca rasa o per una política apresa. Una comparació futura més simètrica permetrà determinar fins a quin punt aquesta conclusió depèn dels límits de les implementacions actuals.
